\section{Diagnóstico del Planeamiento y Unidades de Actuación}

A partir de la base de datos espacial normalizada del Plan General de Ordenación (formato FIP 2.0 \cite{sitcan_planeamiento_2026} integrado en PostGIS mediante el entorno de computación agéntica Google Antigravity \cite{google_antigravity_2026}), se ha analizado la estructura de gestión del suelo en el municipio.

\subsection{Distribución por Categorías del Suelo}

El análisis de las Unidades de Actuación (UAs) delimitadas en el término municipal refleja una concentración determinante en el suelo urbano sujeto a procesos de regularización y renovación:

\begin{table}[htbp]
  \centering
  \caption{Resumen de Unidades de Actuación por Categoría del Suelo}
  \label{tab:resumen_suelo}
  \vspace{0.2cm}
  \begin{tabular}{lrrr}
  \toprule
  \textbf{Categoría del Suelo} & \textbf{Nº UAs} & \textbf{Edificabilidad ($m^2c$)} & \textbf{Superficie ($m^2$)} \\
  \midrule
  Suelo Urbano No Consolidado & 125 & 1.557.793 & 2.915.364 \\
  Suelo Urbanizable Sectorizado & 3 & 80.778 & 112.021 \\
  \midrule
  \textbf{Total} & \textbf{128} & \textbf{1.638.571} & \textbf{3.027.385} \\
  \bottomrule
\end{tabular}

\end{table}

De las 128 unidades delimitadas, el \textbf{97,7\%} (125 UAs) corresponden a \textbf{Suelo Urbano No Consolidado (SUNC)}, sumando una capacidad edificatoria lucrativa de más de 1,55 millones de $m^2c$. Por su parte, el \textbf{Suelo Urbanizable Sectorizado (SUS)} concentra únicamente 3 unidades (ubicadas en el ámbito residencial de El Tablero), con 80.778 $m^2c$.

\begin{figure}[htbp]
  \centering
  \includegraphics[width=0.92\textwidth]{figures/chart_land_classes.pdf}
  \caption{Capacidad edificatoria y número de UAs por categoría de suelo.}
  \label{fig:chart_suelo}
\end{figure}

\subsection{Sistemas de Ejecución y Gestión Urbanística}

El régimen de desarrollo de las actuaciones determina el grado de iniciativa privada frente a la intervención directa de la administración municipal:

\begin{table}[htbp]
  \centering
  \caption{Distribución de Unidades de Actuación según Sistema de Ejecución}
  \label{tab:resumen_sistemas}
  \vspace{0.2cm}
  \begin{tabular}{p{6.8cm}rrr}
  \toprule
  \textbf{Sistema de Ejecución} & \textbf{Nº UAs} & \textbf{\% UAs} & \textbf{Edificabilidad ($m^2c$)} \\
  \midrule
  Sistemas de ejecución privada legalmente establecidos & 75 & 58.6\% & 1.169.062 \\
  Cooperación & 45 & 35.2\% & 463.292 \\
  No especificado & 3 & 2.3\% & 6.217 \\
  Expropiación & 3 & 2.3\% & - \\
  Convenio urbanístico & 2 & 1.6\% & - \\
  \midrule
  \textbf{Total} & \textbf{128} & \textbf{100.0\%} & \textbf{1.638.571} \\
  \bottomrule
\end{tabular}

\end{table}

\begin{itemize}
  \item \textbf{Compensación / Ejecución Privada (58,6\%)}: Representa el mecanismo predominante, acumulando más de 1,16 millones de $m^2c$ edificables donde los propietarios asumen la redacción de los proyectos de urbanización y la cesión de aprovechamientos.
  \item \textbf{Cooperación Pública (35,2\%)}: Abarca 45 unidades (463.292 $m^2c$), donde la ejecución de las obras corresponde al Ayuntamiento a costa de los propietarios.
  \item \textbf{Expropiación y Convenios (3,9\%)}: Reservados para ámbitos con afecciones de sistemas generales o actuaciones asistemáticas singulares.
\end{itemize}
